\paragraph{Terminal extensions of read-once baselines.}
The two closure laws also suffice above every primitive read-once core
operation. The construction uses disjoint supports at each node; it does
not require the source-dependence graph to be acyclic.

\begin{lemma}[Single-mask rigidity]\label{core:lem:read-once-single-mask}
Let $C$ be primitive read-once, and let a terminal mask tuple satisfy
\cref{core:lem:terminal-closure-laws}. If only $\psi_s$ can be nonzero,
then $\psi_s=0$ or $\psi_s=\Phi_s(C)$.
\end{lemma}
\begin{proof}
Every coalition $A$ failing $\Phi_s(C)$ can be reduced to empty by
successively deleting captured inputs, using only foreign-owner capture
contexts disjoint from the current coalition. To prove this, induct on
the tree. At a leaf the failed coalition is already empty. At an
$s$-root both child coalitions fail; run their deletion plans in turn.
At a root $q\ne s$, one child $U$ fails. Run its plan, then fix all
variables of $U$ to $q$. The $q$-root absorbs and captures the entire
other child, deleting what remains.

Captures inside a child lift to the whole word: the queried variable
occurs nowhere outside that child, so the absence of its queried winner
there implies absence in the whole word. A deletion may keep any other
captured inputs: the second closure law and monotonicity allow deleting
an arbitrary subset of the captured set. Thus the inductive plans remain
valid in the whole tree and leave the other child untouched when needed.

Since every foreign mask is zero, the second closure law preserves
$\psi_s$ throughout this deletion. Hence
$\Phi_s(C)(A)=0$ implies $\psi_s(A)=\psi_s(\varnothing)=0$.
Conversely, if $\psi_s\ne0$ and $\Phi_s(C)(A)=1$, then
$c_s(A)$ is the whole support, so the first closure law gives
$\psi_s(A)=1$. This proves the assertion.
\end{proof}

\begin{lemma}[Disjoint union of terminal baselines]
\label{core:lem:terminal-disjoint-union}
Let $C_1,C_2$ be source-separable core operations on disjoint finite
nonempty supports $W_1,W_2$. Suppose the terminal closure laws suffice
for core realizability above each $C_i$. Then they suffice above
$D=C_1+C_2$.
\end{lemma}
\begin{proof}
The owner closures of $D$ are the direct sums
$c_t^D(E)=c_t^1(E\cap W_1)\cup c_t^2(E\cap W_2)$.
For proper component colorings $p,q$, put
\[
 P(p,q)=\bigwedge_t\neg\psi_t(p_t\cup q_t).
\]
Use this formula on all pairs of proper component colorings, even when
their combined palette is full. Only the globally proper pairs will be
needed for the final source predicate.

Fix $q$. If some $\psi_t(q_t)=1$, the section $P(-,q)$ is zero,
the bottom interval predicate. Otherwise its shifted masks
$\psi_t^q(A)=\psi_t(A\cup q_t)$ are normalized and satisfy the
closure laws for $C_1$. For the first law, sandwich
$c_t^1(A)\cup q_t$ between $A\cup q_t$ and
$c_t^D(A\cup q_t)$. For the second, apply the global law to
$E\cup q_t$ and $A\cup q_s$, which are disjoint. Its deletion
also removes $c_t^2(q_t)$ from the fixed $q_s$ block; monotonicity
allows restoring those points, giving exactly the shifted law.
By hypothesis this section is a predicate in $\mathcal I(C_1)$.
Symmetrically every section $P(p,-)$ belongs to $\mathcal I(C_2)$.

These sections give a finite rectangle decomposition. For each $p$ with
nonempty positive row $B_p=\{q:P(p,q)=1\}$, let
\[
 A_p=\bigcap_{q\in B_p}\{u:P(u,q)=1\}.
\]
The interval lattice theorem makes $A_p$ an actual predicate above
$C_1$, while $B_p$ is an actual predicate above $C_2$.
One has $p\in A_p$ and $A_p\times B_p\subseteq P$, so the union
of these rectangles is exactly $P$.
If $F_1,F_2$ realize their factors, the interval operation
\[
 C_2(w_2)+F_1(F_2(X,w_2),w_1)
\]
has baseline $D$ and distinguished predicate equal to their product.
Indeed its bounds are $D$ and $D\cup X$, and at a source probe the
two distinguished coefficients multiply. Union over the rectangles
therefore realizes the required predicate on every globally proper
coloring. Source reconstruction gives the candidate itself.
\end{proof}

\begin{theorem}[Terminal extensions of read-once operations]
\label{core:thm:read-once-terminal-extension}
Let $C$ be a primitive read-once core operation, and give a new source
$X$ normalized separable factors $B_t=\neg\psi_t$. The interval
candidate with these distinguished factors and passive baseline $C$
is core if and only if the two terminal closure laws hold.
The same conclusion holds for disjoint unions of such baselines.
\end{theorem}
\begin{proof}
Necessity is \cref{core:lem:terminal-closure-laws}. For sufficiency
induct on the read-once tree. At a leaf every mask is either zero or
the one-variable singleton test; attaching that variable to $X$ at
the nonzero-mask owners, then adjoining the baseline, realizes the
candidate.

Write a nonleaf baseline as $C=C_1\star_r C_2$ on disjoint supports,
and put $D=C_1+C_2$. The owner closures of $D$ are bounded by those
of $C$. Both terminal laws descend to these smaller closures by
monotonicity: every intermediate value lies between the equal endpoint
values in the law for $C$.

Suppose first $\psi_r\ne0$. The first law implies
$\Phi_r(C)\leq\psi_r$. The preceding lemma and induction give a
core terminal extension $H$ above $D$ with the specified masks.
Then $H\star_r C$ is exactly the target.
Where $C\ne\{r\}$, neither child is $\{r\}$, so $D=C$; also
$H\supseteq C$ prevents an unintended $r$-absorption, and the final
gate leaves $H$ unchanged. Where $C=\{r\}$, the positive $r$-mask
makes the target equal $C$, and the final gate restores that value.

Now suppose $\psi_r=0$. Fixing either whole child to $r$ captures
the entire support. The second law forces every $\psi_t$ to vanish
on coalitions contained in either child alone. For $t\ne r$ and fixed
$A_2\subseteq W_2$, take the pure-$t$ section on the first child.
Its shifted masks satisfy the child laws as in the preceding lemma;
only the $t$-mask can be nonzero, since all the other masks are
evaluated on first-child coalitions alone. Single-mask rigidity makes
this section zero or $\Phi_t(C_1)$. Interchanging the children then
shows that every nonzero global mask is
\[
 \psi_t(A_1\cup A_2)=
 \Phi_t(C_1)(A_1)\Phi_t(C_2)(A_2)=\Phi_t(C)(A_1\cup A_2).
\]
Start a spine at $X$ and attach $C$ at each nonzero-mask owner, then
adjoin $C$. Every gate stays in $C\cup X$, and its distinguished
factors are exactly the prescribed ones. This realizes the candidate.

Finally, repeated application of the disjoint-union lemma proves the
claim for disjoint unions of primitive read-once baselines.
\end{proof}
